\documentclass[a4paper,11pt]{article}
\usepackage{exam-style-341}
\examinehead{341 农业综合知识三 · 程序设计分章练习}{第08章\quad 指针（一）指针基础}
\begin{document}

\unitcover{341 农业综合知识三 · 程序设计分章练习}{第08章\quad 指针（一）指针基础}{练习（共 100 分）}{约 70 分钟}{地址与指针的概念；指针变量的定义、赋值与引用；\code{\&} 与 \code{*} 运算符；指针与数组；指针算术；指针作函数参数；字符串指针；野指针与空指针。}

\begin{knowmap}
\textbf{地址与指针：}内存中每个字节都有编号，称为地址；变量的地址用 \code{\&} 取得；专门用来存放地址的变量称为指针变量。\\[2pt]
\textbf{定义与引用：}\code{int *p;} 定义 p 是指向 int 型变量的指针；\code{p = \&a;} 使 p 指向 a；\code{*p} 表示 p 所指向变量的值，\code{\&} 与 \code{*} 互为逆运算。\\[2pt]
\textbf{指针与数组：}数组名代表数组首元素的地址；\code{p = a;} 后，\code{a[i]}、\code{*(a+i)}、\code{p[i]} 等价；\code{p+1} 不是加 1 个字节，而是跳过 \code{sizeof(所指类型)} 个字节。\\[2pt]
\textbf{指针作参数：}把变量地址传给函数，函数内通过 \code{*p} 修改实参所指变量的值；交换两个变量必须传地址（\code{swap(\&a, \&b)}），只传值无法交换。\\[2pt]
\textbf{字符串指针：}\code{char *s = "hello";} 使 s 指向字符串常量的首字符，可用 \code{puts(s)} 或 \code{\%s} 输出；字符串常量不能通过 s 修改。\\[2pt]
\textbf{常见错误：}指针未初始化（野指针）、空指针解引用、越界访问、指针类型不匹配；使用指针前必须让它指向确定有效的内存单元。
\end{knowmap}

\begin{keyterms}
\textbf{地址}：内存单元的编号，每个字节有唯一地址；变量的地址可用 \code{\&} 取得。\\
\textbf{指针}：一个变量的地址称为该变量的指针；指针是“地址”的形象化说法。\\
\textbf{指针变量}：专门存放地址（指针）的变量，如 \code{int *p;}，其值是一个地址。\\
\textbf{取地址运算符 \code{\&}}：单目运算符，返回操作数（变量）在内存中的地址，如 \code{\&a}。\\
\textbf{间接访问运算符 \code{*}}：单目运算符，作用于指针，返回该指针所指变量的值，如 \code{*p}。\\
\textbf{空指针}：不指向任何有效内存单元的指针，用 \code{NULL} 表示；解引用空指针是非法的。\\
\textbf{指针算术}：指针加、减整数，结果仍是指针，移动的字节数为“整数 $\times$ \code{sizeof(所指类型)}”。\\
\textbf{指针与数组}：数组名是数组首元素的地址，\code{a[i]} 等价于 \code{*(a+i)}，\code{p[i]} 等价于 \code{*(p+i)}。\\
\textbf{指针作函数参数}：实参传地址、形参为指针，函数内通过 \code{*p} 可以修改实参所指的变量。\\
\textbf{字符串指针}：指向字符的指针，通常指向字符串常量的首字符，如 \code{char *s = "China";}。\\
\textbf{野指针}：未初始化或指向已释放内存的指针，其值不确定，解引用会导致未定义行为。\\
\textbf{指针类型}：指针所指数据的类型，如 \code{int *}、\code{char *}；类型决定 \code{*p} 取几个字节以及 \code{p+1} 移动多少字节。
\end{keyterms}

\begin{contrast}{普通变量 / 指针变量（每行依次为：存储内容·占用空间·访问方式）}
\citem{普通变量}{存储数据本身，如 \code{int a = 10;} 存 10}{由数据类型决定，如 int 通常占 4 字节}{直接用变量名访问，如 \code{a}}
\citem{指针变量}{存储另一个变量的地址，如 \code{int *p = \&a;}}{由平台决定，通常占 8 字节（64 位）或 4 字节（32 位）}{先取地址再间接访问，如 \code{*p}}
\end{contrast}

\begin{contrast}{\code{\&a} / \code{*p}（每行依次为：运算类型·含义·结果类型）}
\citem{\code{\&a}}{取地址运算符}{取得变量 a 在内存中的地址}{指针类型，如 \code{int *}}
\citem{\code{*p}}{间接访问运算符}{取得指针 p 所指变量的值}{所指变量的类型，如 \code{int}}
\end{contrast}

\begin{contrast}{数组名 / 指针变量（每行依次为：本质·能否赋值或自增·\code{sizeof} 结果）}
\citem{数组名}{数组首元素的地址，是地址常量}{不能赋值、不能自增；\code{a = p} 和 \code{a++} 都非法}{整个数组的字节数，如 \code{int a[5];} 时为 20}
\citem{指针变量}{存放地址的变量}{可以重新赋值、可以自增或自减}{指针变量本身的大小（通常 8 或 4 字节）}
\end{contrast}

\begin{contrast}{\code{*p++} / \code{(*p)++}（每行依次为：运算顺序·等价写法·对指针和所指值的影响）}
\citem{\code{*p++}}{先取 \code{*p} 的值，再让 p 自增（\code{++} 优先级高于 \code{*}）}{等价于 \code{*(p++)}}{p 后移一个元素，所指变量的值不变}
\citem{\code{(*p)++}}{先取 p 所指变量，再让该变量的值加 1}{等价于 \code{*p = *p + 1}}{p 不动，所指变量的值加 1}
\end{contrast}

\qtypename{一、单项选择题}{10 题，每题 2 分，共 20 分}
\q 若有定义 \code{int a = 10, *p = \&a;}，则 \code{*p} 的值是（\quad）。
\choicesline{A. a 的地址}{B. 10}{C. p 的地址}{D. 0}

\q 运算符 \code{\&} 的作用是（\quad）。
\choicesline{A. 取变量的值}{B. 取变量的地址}{C. 间接访问}{D. 逻辑与运算}

\q 若有定义 \code{int a[5], *p = a;}，则 \code{p + 2} 指向（\quad）。
\choicesline{A. a[1]}{B. a[2]}{C. a[3]}{D. a[2] 的值}

\q 与 \code{a[3]} 等价的是（\quad）。
\choicesline{A. \code{*a + 3}}{B. \code{*(a + 3)}}{C. \code{\&a[3]}}{D. \code{*(a + 3) *}}

\q 指针作函数参数时，实参应传递（\quad）。
\choicesline{A. 变量的值}{B. 变量的地址}{C. 变量名}{D. 变量类型}

\q 执行下列程序段后，\code{*p} 的值是（\quad）。
\begin{lstlisting}[language={[custom]C}]
int a[5] = {10, 20, 30, 40, 50}, *p = a;
p++;
\end{lstlisting}
\choicesline{A. 10}{B. 20}{C. 30}{D. 11}

\q 执行 \code{char *s = "hello";} 后，s 指向（\quad）。
\choicesline{A. 字符 'h'}{B. 整个字符串}{C. 字符 'o'}{D. 空指针}

\q 要使指针 p 不指向任何有效内存单元，应执行（\quad）。
\choicesline{A. \code{p = NULL;}}{B. \code{p = 0.0;}}{C. \code{p = '0';}}{D. \code{p = EOF;}}

\q 若有定义 \code{int a[5], *p = \&a[0];}，则 \code{p + 4} 指向（\quad）。
\choicesline{A. a[3]}{B. a[4]}{C. a[5]}{D. 越界单元}

\q 若有定义 \code{int x = 3, *p = \&x, *q = p;}，则 \code{*q} 的值是（\quad）。
\choicesline{A. 3}{B. p 的地址}{C. q 的地址}{D. 0}

\qtypename{二、多项选择题}{5 题，每题 3 分，共 15 分；每题有两个或两个以上正确项}
\q 关于指针变量，下列说法正确的有（\quad）。
\choicesline{A. 存放的是地址}{B. 定义形式如 \code{int *p;}}{C. \code{*p} 表示所指变量的值}{D. 指针变量本身也占用内存}

\q 关于指针与数组的关系，下列说法正确的有（\quad）。
\choicesline{A. \code{a[i]} 等价于 \code{*(a+i)}}{B. 数组名是数组首元素的地址}
{C. \code{p = a} 使 p 指向 a[0]}{D. 数组名可以像指针一样自增}

\q 关于指针作函数参数，下列说法正确的有（\quad）。
\choicesline{A. 传递的是地址}{B. 可通过 \code{*p} 修改实参所指变量}
{C. 实参可以是变量地址或指针}{D. 一次只能修改一个变量}

\q 关于字符串指针，下列说法正确的有（\quad）。
\choices{A. \code{char *s = "abc";} 使 s 指向字符 'a'}{B. 不能通过 s 修改字符串常量的内容}
{C. 可用 \code{puts(s)} 输出整个字符串}{D. 使用前必须为 s 动态分配内存}

\q 下列做法中属于常见指针错误的有（\quad）。
\choicesline{A. 指针未初始化就解引用}{B. 数组下标越界}
{C. 解引用空指针}{D. 指针指向的内存被释放后继续使用}

\qtypename{三、判断题}{5 题，每题 2 分，共 10 分；正确的画“√”，错误的画“×”并写出改正后的表述}
\q 指针变量本身也占用内存，其大小由平台决定，与所指数据类型无关。\quad（\quad）\quad 改正：\fixblank

\q \code{int *p; *p = 10;} 是正确用法。\quad（\quad）\quad 改正：\fixblank

\q 数组名可以像指针变量一样进行 \code{a++} 运算。\quad（\quad）\quad 改正：\fixblank

\q \code{a[i]} 与 \code{*(a+i)} 等价。\quad（\quad）\quad 改正：\fixblank

\q 指针作函数参数时传递地址，因此可修改实参。\quad（\quad）\quad 改正：\fixblank

\qtypename{四、填空题}{10 空，每空 1 分，共 10 分}
\q 若有定义 \code{int a = 5, *p = \&a;}，则 p 中存放的是\fillblanks[2.0cm]，\code{*p} 的值为\fillblank{}。

\q \code{\&} 是\fillblanks[1.6cm]运算符，\code{*} 是\fillblanks[1.8cm]运算符。

\q 若有定义 \code{int a[5], *p = a;}，则 \code{p + 3} 指向数组元素\fillblanks[1.2cm]。

\q 使指针不指向任何有效内存单元的常量是\fillblanks[1.4cm]。

\q \code{a[i]} 可以写成\fillblanks[1.6cm]。

\q 执行 \code{int a[5]=\{1,2,3,4,5\}, *p=a; p++;} 后，\code{*p} 的值为\fillblank{}。

\q 交换两个变量的值，应把它们的\fillblanks[1.2cm]传给函数。

\q 执行 \code{char *s = "hello";} 后，s 指向字符串的\fillblanks[1.8cm]。

\qtypename{五、程序阅读题}{2 题，每题 5 分，共 10 分}
\q 写出下列程序的运行结果。
\begin{lstlisting}[language={[custom]C}]
#include <stdio.h>
int main(void)
{
    int a[5] = {10, 20, 30, 40, 50};
    int *p = a;
    printf("%d %d %d\n", *p, *(p + 2), p[3]);
    p += 2;
    printf("%d\n", *p);
    return 0;
}
\end{lstlisting}
\anslines{2.4cm}

\q 写出下列程序的运行结果。
\begin{lstlisting}[language={[custom]C}]
#include <stdio.h>
void swap(int *x, int *y)
{
    int t = *x;
    *x = *y;
    *y = t;
}
int main(void)
{
    int a = 3, b = 8;
    int *p = &a, *q = &b;
    swap(p, q);
    printf("%d %d\n", a, b);
    return 0;
}
\end{lstlisting}
\anslines{2.6cm}

\qtypename{六、简答题}{3 题，每题 5 分，共 15 分}
\q 什么是指针？指针变量与普通变量有何区别？\code{\&} 和 \code{*} 两个运算符各起什么作用？\anslines{3.2cm}

\q 指针与数组有何关系？为什么 \code{a[i]} 等价于 \code{*(a+i)}？数组名与指针变量有何不同？\anslines{3.4cm}

\q 指针作函数参数有什么优点？请举例说明如何通过指针在函数中交换两个变量的值。\anslines{3.2cm}

\newpage
\qtypename{七、综合题}{1 题，共 20 分}
\q 编写一个完整的 C 程序，要求：
\par\vspace{0.15em}
\noindent（1）定义函数 \code{void find(int a[], int n, int *max, int *min)}，通过指针参数“带回”数组中 $n$ 个元素的最大值和最小值；\\
\noindent（2）在 \code{main} 中读入 $n$ 及 $n$ 个整数，调用 \code{find} 后输出最大值和最小值。

\noindent 输入示例：\code{6  4 -2 9 7 0 3}，输出应为 \code{max=9,min=-2}。请写出完整程序。\anslines{4.0cm}

\end{document}
